\documentclass{ximera}
% xmPreamble.tex en xmPrintstyle.sty worden automatisch geladen door ximera.cls (v2.7.8);
% NIET nog eens \input{./preamble.tex} of \addPrintStyle{.} doen (geeft 'already defined'-fouten).

\begin{document}
\author{Alexander Holvoet}
\xmtitle{Een hoop zand op een vierkant}{}

\begin{exercise}
    Wat denk je te krijgen als je zand giet op een vierkant plaatje? Doe een voorspelling
    en voer daarna pas het experiment uit.
    \begin{oplossing}
        Je krijgt een rechte piramide met het vierkant als grondvlak.
    \end{oplossing}
\end{exercise}

% TODO (afbeelding ontbreekt): in de bron stond hier \afbeelding[5cm]{LPiramide}
% (label fig:piramide). Lijntekening niet aanwezig in img/.

\begin{exercise}
    Meet opnieuw de rusthoek en vergelijk met de rusthoek bij de kegel.
    \begin{oplossing}
        De rusthoek is dezelfde als bij de kegel, tenminste als je hetzelfde materiaal
        gebruikt hebt.
    \end{oplossing}
\end{exercise}

\begin{exercise}
    Teken het bovenaanzicht van de zandhopen met een rond grondvlak en met een vierkant
    grondvlak.
    \begin{hint}
        Wat gebeurt er met de top van de kegel in bovenaanzicht? En met de ribben van de
        piramide?
    \end{hint}
    \begin{oplossing}
        Bij het ronde grondvlak krijg je een cirkel met een puntje. Bij het vierkante
        grondvlak krijg je een vierkant met diagonalen. Het puntje staat voor de top van
        de kegel, de diagonalen staan voor de ribben van de piramide. Uit het bovenaanzicht
        kun je niet afleiden hoe groot de rusthoek is.
    \end{oplossing}
\end{exercise}

% TODO (afbeelding ontbreekt): in de bron stond hier \afbeelding[7cm]{Lbovenzichtkegel}
% (label fig:bovenzicht). Lijntekening niet aanwezig in img/.

\end{document}
